% !TEX root = main.tex
% scPhyTr --- ICML-style methods paper (DRAFT scaffold).
% Self-contained preamble so it compiles with `tectonic main.tex` out of the box.
% To submit to ICML: replace the geometry/font block below with the official
% \usepackage{icml2025} style file and \twocolumn layout.
\documentclass[11pt]{article}

\usepackage[margin=1in]{geometry}
\usepackage{amsmath,amssymb,amsthm}
\usepackage{bm}
\usepackage{booktabs}
\usepackage{algorithm}
\usepackage{algpseudocode}
\usepackage{microtype}
\usepackage{natbib}
\usepackage{xcolor}
\usepackage[colorlinks=true,allcolors=blue]{hyperref}
\usepackage[capitalize]{cleveref}

\newtheorem{proposition}{Proposition}
\newtheorem{lemma}{Lemma}
\newtheorem{definition}{Definition}

\newcommand{\scphytr}{\textsc{scPhyTr}}
\newcommand{\RR}{\mathbb{R}}
\newcommand{\E}{\mathbb{E}}
\newcommand{\N}{\mathcal{N}}
\newcommand{\KK}{K}
\newcommand{\bigO}{\mathcal{O}}

\title{\bf Exact, Linear-Time Inference for Latent Gaussian Evolutionary
Models on Single-Cell Phylogenies}
\author{\textit{(author list)}}
\date{Draft --- \today}

\begin{document}
\maketitle

\begin{abstract}
Single-cell lineage tracing now reconstructs a phylogeny over the individual cells of a
growing tissue or tumour and profiles each cell's transcriptome, raising a statistical
question of how gene expression \emph{evolves} along a cellular genealogy: how heritable a
program is, whether it is under selection, which genes co-evolve, and where the tempo of
evolution shifts. We cast this as inference in a \emph{latent Gaussian model on a tree}: a
Brownian-motion or Ornstein--Uhlenbeck process places a latent state at every node, observed
at the leaves through an explicit (Poisson / negative-binomial) count likelihood. The joint
law is a Gaussian Markov random field whose precision inherits the tree's sparsity, which we
exploit for a \emph{sparsity-preserving Laplace} inference that evaluates the marginal
likelihood and its exact gradients in time \emph{linear} in the number of cells---exact when
the observation is Gaussian, and tight for log-concave count decoders---scaling to the gene
dimension through a low-rank phylogenetic factor model in $\bigO(nk^3)$. This yields an
exact-marginal, maximum-likelihood alternative to the amortized variational tree models and
to the dense $\bigO(n^3)$ likelihoods of classical phylogenetic comparative methods, and a
calibrated penalized-likelihood search for evolutionary \emph{rate shifts} that matches the
discoveries of reversible-jump MCMC roughly two orders of magnitude faster. On ground-truth
simulations and on real lineage-tracing data, \scphytr{} recovers heritability, deconfounded
gene--gene co-evolution, and clade-specific rate shifts where autocorrelation- and
graph-based baselines are confounded by shared ancestry.
\end{abstract}

\section{Introduction}
\label{sec:intro}

CRISPR-based lineage tracing records, in an evolving population of cells, a heritable barcode
that is read out alongside single-cell RNA sequencing, so that one obtains both a
reconstructed \emph{phylogeny} over thousands of cells and each cell's expression profile.
This makes it possible, for the first time at single-cell resolution, to ask the questions
that phylogenetic comparative methods (PCMs) ask of species: is a trait \emph{heritable}
along the genealogy or rapidly \emph{plastic}; is it under \emph{stabilizing or directional
selection}; which traits \emph{co-evolve}; and does the \emph{rate} of evolution accelerate
in particular subclones? Existing tools answer one of these at a time and with mismatched
assumptions---phylogenetic autocorrelation for heritability, dense Gaussian-on-log PCM for
selection, graph autocorrelation for modules, reversible-jump MCMC for rate shifts---and most
treat normalized expression as a Gaussian trait, which fails at the sparse depths typical of
single-cell data.

We give a single generative model and an inference engine that produce all of these read-outs
from one fit. The model is a latent BM/OU process on the tree observed through a count
likelihood (\Cref{sec:model}); the engine exploits the tree-induced sparsity of the latent
precision to compute the marginal likelihood and its gradients exactly and in linear time
(\Cref{sec:inference}). Our contributions are:
\begin{itemize}
  \item \textbf{Linear-time exact inference for latent Gaussian tree models with non-conjugate
  observations.} A sparsity-preserving Laplace approximation evaluates the marginal in
  $\bigO(np^3)$ ($\bigO(n)$ univariate), exact when the decoder is Gaussian, with exact
  gradients and an EM whose E-step is a Rauch--Tung--Striebel tree smoother
  (\Cref{sec:inference}). This is an exact-MLE counterpart to amortized variational tree
  models and to dense $\bigO(n^3)$ PCM.
  \item \textbf{Scalability to the transcriptome} through a low-rank Poisson phylogenetic
  factor model, $K=WW^\top$, fit by Laplace-EM in $\bigO(nk^3)$, recovering the deconfounded
  gene--gene evolutionary correlation matrix without forming the $p\times p$ object
  (\Cref{sec:estimation}).
  \item \textbf{A calibrated penalized-likelihood search for clade-specific rate shifts}
  (\Cref{sec:estimation}), the maximum-penalized-likelihood counterpart to RevBayes'
  reversible-jump MCMC, with a selection-bias correction that controls the null false-positive
  rate.
  \item \textbf{Empirical evidence} (\Cref{sec:experiments}) that, on ground-truth
  simulations and real lineage-tracing data, these read-outs are recovered where PATH,
  Hotspot, EvoGeneX and RevBayes are confounded, mis-calibrated, or orders of magnitude
  slower.
\end{itemize}

\section{Model: a latent Gaussian process on the tree}
\label{sec:model}

Let $T$ be a rooted tree with nodes $u$ and leaf set $\mathcal L$ (the profiled cells). For
$p$ genes, a latent state $z_u\in\RR^p$ (interpreted as log-expression) evolves along edges.
Under \emph{Brownian motion} with diffusion matrix $K\succ 0$, an edge of length $t$ adds an
independent Gaussian increment,
\begin{equation}
z_{\text{child}}\mid z_{\text{parent}} \sim \N\!\big(z_{\text{parent}},\; t\,K\big),
\label{eq:bm}
\end{equation}
so the diagonal of $K$ is the per-gene rate and its off-diagonal the gene--gene evolutionary
covariance. The \emph{Ornstein--Uhlenbeck} process adds a pull toward an optimum $\theta$,
$\mathrm dz=-\alpha(z-\theta)\,\mathrm ds+K^{1/2}\mathrm dW$, integrating over an edge to a
linear-Gaussian transition with contraction $\phi=e^{-\alpha t}$ and variance factor
$v(t)=(1-e^{-2\alpha t})/(2\alpha)$; clade-specific optima or rates are encoded by painting
\emph{regimes} on the tree.

\paragraph{Observation model.} The leaf latent is observed through a likelihood that
factorizes across entries given $z$,
\begin{equation}
Y_{ig}\mid z_{ig} \sim p(\cdot\mid z_{ig}),\qquad
i\in\mathcal L,\; g\in[p],
\end{equation}
with \scphytr{} implementing a Poisson model $Y\sim\mathrm{Poisson}(S\,e^{z})$ with per-cell
size factors $S$, a negative-binomial extension for within-clone overdispersion, and a
Gaussian model for directly measured traits. Crucially, \emph{all evolutionary parameters live
in the latent prior and none in the observation model}, so any observation type reuses the
same inference engine and the estimated $K$ is a property of evolution, not of measurement
noise.

\paragraph{The joint law is a GMRF.} Stacking the latents of all $N$ nodes, the prior is
$\N(Z;m,Q^{-1})$ with precision $Q=A\otimes K^{-1}$, where $A$ is the tree's (sparse)
structure matrix---nonzero only on the diagonal and parent--child pairs. The posterior is
\begin{equation}
p(Z\mid Y)\;\propto\;\exp\!\Big\{\ell(F)-\tfrac12 (Z-m)^\top Q (Z-m)\Big\},
\label{eq:posterior}
\end{equation}
with $F=Z_{\mathcal L}$ and $\ell$ the decoder log-likelihood. Two facts drive everything
below: $\ell$ is concave for log-link Poisson/Gaussian (so \eqref{eq:posterior} is unimodal),
and the posterior precision $Q+W$ with $W=-\nabla^2\ell$ block-diagonal over leaves
\emph{preserves the tree sparsity}.

\section{Linear-time exact inference}
\label{sec:inference}

\subsection{Gaussian observations: exact pruning}
When the decoder is Gaussian the latents can be integrated out in closed form. Rather than
forming the dense $N\!\times\!N$ covariance (an $\bigO(n^3)$ operation), \scphytr{} evaluates
the marginal by Felsenstein's pruning---a single post-order pass fusing children---which is
exactly Gaussian elimination on the sparse precision \eqref{eq:posterior}.

\begin{proposition}[Exactness and cost of pruning]
\label{prop:pruning}
For a Gaussian (or directly observed) trait, the pruning recursion returns the exact marginal
log-likelihood $\log p(Y\mid\eta)$ of the BM/OU model in $\bigO(n)$ time for a single trait
and $\bigO(np^3)$ for $p$ correlated traits, where the $p^3$ is one Cholesky factorization of
$K$ per node and introduces no fill-in.
\end{proposition}

\subsection{Non-conjugate observations: a sparsity-preserving Laplace approximation}
For count decoders the posterior \eqref{eq:posterior} is non-Gaussian but log-concave. We
approximate it by the Gaussian matching its mode and curvature. The mode $\hat Z=\arg\min\Psi$,
$\Psi(Z)=-\ell(F)+\tfrac12(Z-m)^\top Q(Z-m)$, is found by Newton's method, each step solving
\begin{equation}
(Q+W)\,\delta = -\nabla\Psi,\qquad \nabla\Psi=Q(Z-m)-\nabla\ell,
\end{equation}
by tree (block) Gaussian elimination, and the Laplace marginal is the Gaussian-at-the-mode
value.

\begin{proposition}[Sparsity preservation and cost]
\label{prop:laplace}
Because $W=-\nabla^2\ell$ is block-diagonal over leaves, the Hessian $H=Q+W$ has the same
tree sparsity as $Q$; its Cholesky factorization in leaves-to-root order introduces no
fill-in, so the Laplace marginal, its gradient, and the posterior edge moments (via a
Rauch--Tung--Striebel smoother reusing the elimination factors) all cost $\bigO(np^3)$
($\bigO(n)$ univariate). The approximation is exact when $\ell$ is Gaussian and its error is
governed by the third derivative of $\ell$ at the mode.
\end{proposition}

The Laplace step is the default E-step; importance sampling and Hamiltonian Monte Carlo serve
as certificates and fallbacks when the posterior is far from Gaussian (\Cref{app:engines}).

\subsection{Scaling to many genes: a low-rank factor decoder}
The full $K$ is $\bigO(p^2)$ and the dense estimator $\bigO(p^3)$. \scphytr{} parameterizes
$K$ in low rank: $k$ latent factors $x_u\in\RR^k$ evolve as independent BM on the tree and
load onto genes, $\eta_{i}=\mu+Wx_i$, $Y_i\sim\mathrm{Poisson}(S_i e^{\eta_i})$, so
$K=WW^\top$. A single factor influences every gene, so the per-leaf curvature is a full
$k\times k$ block $W^\top\!\operatorname{diag}(S_ie^{\eta_i})W\succeq 0$ rather than diagonal;
this is the \emph{only} change the engine needs, and it leaves \Cref{prop:laplace} intact with
cost $\bigO(nk^3)$.

\section{Estimation and structure search}
\label{sec:estimation}

With the latents integrated out, evolutionary parameters are estimated by maximizing the
(exact or Laplace) marginal: closed form for $K$ when traits are directly observed, otherwise
direct optimization over $(\alpha,\sigma^2,\theta)$, or Expectation--Maximization when extra
latents (factor scores, NB dispersion) are present---an RTS-smoother E-step and a JAX-gradient
M-step. The shared engine specializes per task; one read-out is the exception.

\paragraph{Heritability (Pagel's $\lambda$): the dense exception.} Phylogenetic signal is read
by Pagel's $\lambda$, scaling the off-diagonals of the tip covariance by $\lambda\in[0,1]$
($\lambda{=}1$ full tree, $\lambda{=}0$ i.i.d.\ tips). Scaling \emph{only} the off-diagonals
yields a matrix that is no longer the covariance of a BM on the tree, so it loses the
tree-Markov sparsity and is estimated by a dense $\bigO(n^3)$ factorization---the lone
read-out that leaves the linear-time regime.

\paragraph{Clade-specific rate shifts.} For multi-rate BM (the diffusion rate $\sigma^2$ shifts
between painted regimes), given regimes the per-regime rates are fit by the same linear-time
pruning of \Cref{prop:pruning}. When the shifts are unknown they are found \emph{de novo} by
greedy forward selection: add the branch whose rate shift most improves a penalized criterion,
each candidate scored by its $\bigO(n)$ marginal.

\begin{proposition}[Calibration of the de-novo search]
\label{prop:calibration}
Selecting the best of $m$ candidate branches biases the maximum BIC improvement upward; adding
a location penalty of $2\log m$ per shift corrects this selection bias. Empirically it reduces
the false-positive rate on homogeneous-rate null trees from $50\%$ to $0\%$ while retaining
power (\Cref{sec:experiments}).
\end{proposition}

This is the maximum-penalized-likelihood counterpart to the reversible-jump MCMC of RevBayes:
it returns one configuration in seconds rather than a posterior in minutes.

\section{Experiments}
\label{sec:experiments}

We compare \scphytr{} to PATH (Moran's $I$ autocorrelation), Hotspot (graph autocorrelation),
EvoGeneX (replicate-aware Gaussian OU), and RevBayes (reversible-jump rate shifts), on
ground-truth simulations and on real KP-Tracer and melanoma-subline lineage-tracing data.
\Cref{tab:results} summarizes the headline results; full protocols are in \Cref{app:experiments}.

\begin{table}[t]
\centering
\small
\begin{tabular}{@{}lll@{}}
\toprule
Task / regime & \scphytr{} & baseline \\
\midrule
Heritability AUROC, depth $1$ & $0.82$ & EvoGeneX $0.66$, PATH $0.65$ \\
Co-evolution false-positive rate & $0.4\%$ & PATH $96\%$, Hotspot $73\%$ \\
PFA module recovery, $p{=}8000$ & AUROC $1.0$, ${\sim}15$s & dense $K$ infeasible \\
Rate shift localize ($16\times$) & $100\%$ & RevBayes $75\%$ \\
Rate magnitude ($16\times$ true) & $17\times$ & RevBayes $3.6\times$ \\
Rate-shift runtime / dataset & ${\sim}1$s & RevBayes ${\sim}85$s \\
Null rate-shift FPR & $0\%$ & --- \\
\bottomrule
\end{tabular}
\caption{Headline benchmark results. Simulations use planted ground truth; depth is mean
counts/cell. Source scripts in \texttt{analysis/benchmark/}.}
\label{tab:results}
\end{table}

\paragraph{Heritability and the count model.} On simulations sweeping sequencing depth, the
count model retains power to detect heritable expression where Gaussian-on-log methods lose it
(AUROC $0.82$ vs $0.65$--$0.66$ at depth $1$; all converge at deep coverage). On real
lineage-tracing data \scphytr{}'s $\lambda$ reproduces PATH's Moran's $I$ ranking.

\paragraph{Co-evolution without tree confounding.} With a planted low-rank $K$ on real tumour
trees, \scphytr{}'s deconfounded estimate holds the false co-evolution rate among
evolutionarily independent genes at $0.4\%$ at full power, while PATH's cross-correlation
($96\%$) and Hotspot ($73\%$) manufacture co-evolution from shared ancestry. The factor model
recovers the planted module to $p=8000$ genes in seconds where the dense $p\times p$ estimate
is rank-deficient and infeasible.

\paragraph{Rate shifts vs reversible-jump MCMC.} On trees with a known rate shift, \scphytr{}
localizes large shifts more often than RevBayes' highest-posterior branch and recovers the
rate magnitude near-unbiased, whereas the reversible-jump posterior shrinks it toward the null;
\scphytr{} runs ${\sim}100\times$ faster. The de-novo search is calibrated
(\Cref{prop:calibration}).

\section{Related work}
\label{sec:related}

\scphytr{} sits between three camps. \emph{Neural tree models} (TreeVAE, LORACs, VOUS) place a
Gaussian random walk on a (often latent) tree and decode with a neural network trained by
amortized variational inference; we keep a log-linear decoder in exchange for an exact marginal,
exact gradients, and interpretable evolutionary parameters. \emph{Phylogenetic comparative
methods} (EvoGeneX, CAGEE) fit the same BM/OU models but to bulk/pseudobulk expression, one
gene at a time, with dense $\bigO(n^3)$ likelihoods; \scphytr{} lifts this to single-cell
resolution, count-based and multivariate, in linear time. \emph{Moment- and graph-statistics}
(PATH, Hotspot) summarize the tree rather than model it, and are confounded by shared ancestry
in exactly the co-evolution regime our deconfounded $K$ targets. \emph{RevBayes} is the
Bayesian end of the PCM axis, fitting rate shifts by reversible-jump MCMC; \scphytr{} trades the
posterior for a calibrated, linear-time penalized-likelihood estimate. Connections to INLA,
Gaussian-process classification, and Kalman/RTS smoothing are detailed in \Cref{app:related}.

\section{Conclusion}
\scphytr{} casts the study of expression evolution on single-cell phylogenies as inference in a
latent Gaussian tree model with a count decoder, and shows that the tree's sparsity yields an
exact, linear-time marginal likelihood that subsumes the read-outs of specialized tools in one
calibrated, interpretable model. Limitations and extensions---small-tree regularization of
$\lambda$, a count-model de-novo search, and proliferation dynamics---are discussed in
\Cref{app:limitations}.

\appendix
\section{Inference engines: importance sampling and HMC}
\label{app:engines}
\emph{(To import from \texttt{docs/02\_inference\_engines.md}: the FFBS simulation smoother,
the self-normalized importance-sampling weights and marginal estimator, and the NUTS target;
plus the head-to-head agreement of the three engines.)}

\section{Experimental protocols}
\label{app:experiments}
\emph{(Datasets, simulation generators, baseline configurations, and full tables; source in
\texttt{analysis/benchmark/}.)}

\section{Related work, extended}
\label{app:related}
\emph{(INLA / latent Gaussian models; GP classification Laplace; Kalman/RTS; PATH/PATHpro,
EvoGeneX/CAGEE, Hotspot/PhyloVision, TreeVAE/LORACs/VOUS, RevBayes.)}

\section{Limitations}
\label{app:limitations}
\emph{(Laplace error for very sparse genes; small-tree $\lambda$ overfitting and a bootstrap
fix; de-novo rate shifts currently on the Gaussian path; tree taken as given.)}

\bibliographystyle{plainnat}
\bibliography{../manuscript/references}

\end{document}
